\documentclass[../../../include/open-logic-section]{subfiles}

\begin{document}

\olfileid[pt]{sth}{ordinals}{replacement}
\olsection{Substituição}

Na \olref[sth][ordinals][vn]{sec}, motivamos a introdução dos ordinais ao
sugerir que os poderíamos tratar como tipos de ordem, isto é, como
representantes canónicos de boas ordenações. Para que isso funcione,
seria preciso provar que \emph{toda a boa ordenação é isomorfa a algum
ordinal}. Isso permitiria definir $\ordtype{A, <}$ como o ordinal
$\alpha$ tal que $\tuple{A, <} \isomorphic \alpha$.

Infelizmente, \emph{não} podemos provar o resultado desejado apenas
com os axiomas introduzidos até aqui. (Veremos por que em
\olref[replacement][strength]{sec}, mas por agora basta: não podemos.)
Precisamos de uma ideia nova, e eis-a:

\begin{axiom}[Esquema de Substituição]
Para qualquer fórmula $\phi(x, y)$, o seguinte é um axioma:
\begin{quote}
	para qualquer $A$, se $(\forall x \in A)\lexists![y][\phi(x,y)]$, então
	$\Setabs{y}{(\exists x \in A)\phi(x,y)}$ existe.
\end{quote}
\end{axiom}
\noindent
Como na Separação, isto é um esquema: produz infinitamente muitos
axiomas, um para cada uma das infinitas $\phi$ distintas. E pode
igualmente ser (e normalmente é) escrito assim:

\begin{defish}
Para qualquer fórmula $\phi(x,y)$ que não contenha ``$B$'', o seguinte é um axioma:
\[
\forall A[(\forall x \in A)\lexists![y][\phi(x,y)] \lif \exists B\forall y (y \in B \liff (\exists x \in A)\phi(x,y))]
\]
\end{defish}

À primeira vista, contudo, esta é uma fórmula bastante intrincada. A
consequência imediata seguinte da Substituição dá provavelmente uma
expressão \emph{mais clara} à ideia intuitiva com que trabalhamos:\footnote{Um termo é uma expressão que designa exatamente um objeto (em qualquer instanciação das suas variáveis livres). Por exemplo, ``$\emptyset$'' é um termo que designa o conjunto vazio; ``$\{x\}$'' é um termo que designa o singular de $x$ (seja qual for $x$); ``$x \cup y$'' é um termo que designa a união de $x$ e de $y$ (sejam quais forem).}

\begin{cor}
Para qualquer termo $\tau(x)$, e qualquer conjunto $A$, este conjunto existe:
\[
	\Setabs{\tau(x)}{x \in A} = \Setabs{y}{(\exists x \in A)y = \tau(x)}.
\]
\end{cor}

\begin{proof}
Como $\tau$ é um \emph{termo}, $\forall x \lexists![y][\tau(x) = y]$.
\emph{A fortiori}, $(\forall x \in A)\lexists![y][\tau(x) = y]$. Logo
$\Setabs{y}{(\exists x \in A)\tau(x) = y}$ existe pela Substituição.
\end{proof}
\noindent
Isto sugere que ``Substituição'' é um bom nome para o axioma: dado um
conjunto $A$, pode formar um novo conjunto, $\Setabs{\tau(x)}{x \in A}$,
substituindo cada membro de $A$ pela sua imagem sob~$\tau$. De fato,
seguindo a notação para a imagem de um conjunto sob uma função,
poderíamos escrever $\funimage{\tau}{A}$ para $\Setabs{\tau(x)}{x \in A}$.

O essencial, porém, é que $\tau$ é um \emph{termo}. Não tem de ser (um
nome de) uma \emph{função}, no sentido de \olref[sfr][fun][rel]{sec},
isto é, um certo conjunto de pares ordenados. Afinal, se $f$ é uma
função (nesse sentido), então o conjunto $\funimage{f}{A} = \Setabs{f(x)}{x \in
A}$ é apenas um subconjunto particular de $\ran{f}$, e isso já está
garantido a existir apenas com os axiomas de~$\Zminus$.\footnote{Baste
considerar $\Setabs{y \in \bigcup \bigcup f}{(\exists x \in A)y =
f(x)}$.} A Substituição, em contraste, é um acréscimo \emph{poderoso} aos
nossos axiomas, como veremos no \olref[sth][replacement][]{chap}.


\end{document}
